% The web pair's lineage figure, same conventions as fig:paths: role names
% not project names, {braces} = a step the reader cannot follow by the name
% at the call site. Drop into the v3 frontend section.
\begin{figure}[H]
\footnotesize
\begin{verbatim}
plain JS             React
1.1k, all in page    1.1k + 283k runtime

event listener       {synthetic event}
  |                    |
handler              {priority lane}
  |                    |
filter + sort        state setter
  |                    |
template string      {hook state store}
  |                    |
innerHTML write      component function
  |                    |
browser paints       filter + sort
                       |
                     {reconciler diff}
                       |
                     {key identity}
                       |
                     {commit phase}
                       |
                     {prop application}
                       |
                     browser paints

  |   a step to a name at the call site,
      resolvable by one search
 {}  a step resolvable only in the
     framework's source or conventions
\end{verbatim}
\caption{The third pair, the same interaction on the web: a keystroke in
the filter until the browser paints the rows. The browser engine (parse,
layout, paint) is floor on both sides by the rule that makes MySQL floor
in the Java pair. Every plain step is in the page's own 1,092 tokens; the
braced steps resolve in the 282,717 tokens of the runtime's source
(react, react-dom, the reconciler, the scheduler, at the measured
version), or, for the delegated event's target and the controlled
input's value tracking, in version convention. The two sides' own code
is the same size (1,092 against 1,145 tokens, matched behavior,
identical rendered DOM); the difference is the machinery behind the
braces.}
\label{fig:webpaths}
\end{figure}
